% =============================================================
%  Kisi-Kisi Ujian Tengah Semester (UTS)
%  IF25-40304 · Pembelajaran Mesin Multimodal · ITERA
%  Cakupan: Pertemuan 01--07
% =============================================================
\documentclass[12pt, a4paper]{article}

\usepackage{tugas-style}

% ── Identitas dokumen ───────────────────────────────────────
\renewcommand{\doctitle}{Kisi-Kisi Ujian Tengah Semester}
\renewcommand{\duedate}{\color{ctpBlue}\faCalendarCheck\ Ujian Tengah Semester (Pertemuan 08) \textbar\ Cakupan Pertemuan 01--07 \faCalendarCheck}

\lhead{{\footnotesize\color{ctpOverlay}\textsc{Kisi-Kisi UTS}}}
\rhead{{\footnotesize\color{ctpOverlay}\textsc{\coursenumber\ --\ \coursetitle}}}

\begin{document}

% ── HEADER ──────────────────────────────────────────────────
\makeassignmentheader

\hrule width0.3\textwidth
\revisioninfo{%
  \vhEntry{0.1}{25 September 2026}{I.W.W.}{Draf inisial kisi-kisi UTS.}
}
\hrule
\vskip .3in

\tableofcontents
\clearpage

% =============================================================
\section{Informasi Umum}\label{sec:info}

Dokumen ini memuat kisi-kisi Ujian Tengah Semester (UTS) mata kuliah
\textbf{Pembelajaran Mesin Multimodal}. Kisi-kisi disusun dari capaian
pembelajaran dan materi Pertemuan 01--07, dan menjadi acuan baik bagi dosen
dalam menyusun soal maupun bagi mahasiswa dalam mempersiapkan diri.

\renewcommand{\arraystretch}{1.35}
\begin{center}\small
\begin{tabular}{|L{4.2cm}|L{10.2cm}|}
  \hline
  \thead{\color{ctpBlue} Aspek} & \thead{\color{ctpBlue}Keterangan} \\
  \hline
  Mata kuliah & Pembelajaran Mesin Multimodal (IF25-40304) --- 3 SKS \\
  \hline
  Semester & Ganjil 2026/2027 \\
  \hline
  Bentuk ujian & Ujian tertulis \textbf{individu}, \textbf{tutup buku} \\
  \hline
  Alat bantu & Kalkulator sederhana. Tanpa internet, tanpa perangkat komputasi. \\
  \hline
  Waktu & \textbf{100 menit} \\
  \hline
  Cakupan & Pertemuan 01--07 (paruh pertama semester) \\
  \hline
  Jumlah butir & 6 butir: 4 esai konseptual + 2 analisis/studi kasus \\
  \hline
  Total poin & 100 poin \\
  \hline
  Bobot nilai akhir & \textbf{25\%} \\
  \hline
\end{tabular}
\end{center}

% =============================================================
\section{Capaian Pembelajaran yang Diuji}\label{sec:cpmk}

UTS menguji lima capaian pembelajaran mata kuliah (CPMK) berikut. Kode level
kognitif mengikuti taksonomi Bloom (C2 = memahami, C3 = menerapkan,
C4 = menganalisis, C5 = mengevaluasi).

\begin{objectives}
  \item \textbf{CPMK-1} (C2) --- Menjelaskan konsep dasar pembelajaran mesin
        multimodal dan tiga pilar utamanya --- \emph{representasi},
        \emph{penyejajaran}, dan \emph{fusi} --- beserta tantangannya.
  \item \textbf{CPMK-2} (C3--C4) --- Menganalisis dan memilih representasi yang
        tepat untuk tiap modalitas (teks, gambar, audio/video).
  \item \textbf{CPMK-3} (C4) --- Membandingkan dan menerapkan strategi fusi
        multimodal beserta \emph{trade-off}-nya.
  \item \textbf{CPMK-4} (C2--C4) --- Menjelaskan dan menganalisis mekanisme
        penyejajaran (\emph{temporal} dan \emph{struktural}) serta peran
        \emph{attention} sebagai penyejajar.
  \item \textbf{CPMK-5} (C5) --- Mengevaluasi rancangan arsitektur multimodal
        pada suatu studi kasus dan mengusulkan perbaikan berbasis alasan.
\end{objectives}

% =============================================================
\section{Cakupan Materi}\label{sec:cakupan}

\renewcommand{\arraystretch}{1.3}
\begin{center}\small
\begin{tabular}{|C{1.4cm}|L{5.4cm}|L{7.4cm}|}
  \hline
  \thead{\color{ctpBlue}W} & \thead{\color{ctpBlue}Topik} & \thead{\color{ctpBlue}Butir Esensial yang Diuji} \\
  \hline
  01 & Pengantar PM Multimodal & Definisi dan motivasi multimodal; tiga pilar (representasi, penyejajaran, fusi) dan tantangannya. \\
  \hline
  02 & Representasi Teks & \emph{Word embeddings}, \emph{contextual embeddings} (BERT), RNN/LSTM, \emph{transformer encoder}. \\
  \hline
  03 & Representasi Gambar & CNN (ResNet) untuk fitur hierarkis; Vision Transformer (ViT) berbasis \emph{patch}. \\
  \hline
  04 & Representasi Audio \& Video & Spektrogram, MFCC, 3D CNN, RNN/LSTM untuk sekuens. \\
  \hline
  05 & Fusi I: Early \& Late & Titik penggabungan, kelebihan, dan kelemahan masing-masing. \\
  \hline
  06 & Fusi II: Intermediate \& Hybrid & Penggabungan fitur terenkode; kombinasi beberapa titik; \emph{trade-off}. \\
  \hline
  07 & Penyejajaran (Alignment) & Penyejajaran temporal dan struktural; attention sebagai penyejajar adaptif. \\
  \hline
\end{tabular}
\end{center}

% =============================================================
\section{Kisi-Kisi Soal}\label{sec:kisi}

\renewcommand{\arraystretch}{1.35}
\setlength{\tabcolsep}{4pt}
\begin{center}\footnotesize
\begin{tabular}{|C{0.7cm}|C{1.0cm}|C{1.2cm}|L{2.1cm}|L{4.3cm}|C{1.0cm}|C{1.7cm}|C{0.8cm}|}
  \hline
  \thead{\color{ctpBlue}No} & \thead{\color{ctpBlue}Soal} & \thead{\color{ctpBlue}CPMK} & \thead{\color{ctpBlue}Materi} & \thead{\color{ctpBlue}Indikator} & \thead{\color{ctpBlue}Level} & \thead{\color{ctpBlue}Bentuk} & \thead{\color{ctpBlue}Poin} \\
  \hline
  1 & A1 & 1 & W01 & Menjelaskan tiga pilar beserta contoh dan risiko bila lemah. & C2 & Esai & 15 \\
  \hline
  2 & A2 & 2 & W02--04 & Membandingkan representasi teks, gambar, dan audio/video. & C3 & Esai & 15 \\
  \hline
  3 & A3 & 3 & W05--06 & Membandingkan empat strategi fusi dan memberi contoh penerapan. & C4 & Esai & 15 \\
  \hline
  4 & A4 & 4 & W07 & Menjelaskan penyejajaran temporal/struktural dan peran attention. & C2--C3 & Esai & 15 \\
  \hline
  5 & B1 & 2, 4 & W02--04, W07 & Menganalisis diagram arsitektur dan mengusulkan perbaikan. & C4 & Analisis & 20 \\
  \hline
  6 & B2 & 3, 5 & W04--07 & Merancang pipeline multimodal dan menimbang \emph{trade-off}. & C5 & Studi kasus & 20 \\
  \hline
  \multicolumn{7}{|r|}{\bfseries Total} & \bfseries 100 \\
  \hline
\end{tabular}
\end{center}
\setlength{\tabcolsep}{6pt}

% =============================================================
\section{Distribusi Level Kognitif dan Estimasi Waktu}\label{sec:distribusi}

\renewcommand{\arraystretch}{1.3}
\begin{center}\small
\begin{tabular}{|L{3.2cm}|C{2.2cm}|C{2.6cm}|C{2.6cm}|}
  \hline
  \thead{\color{ctpBlue}Level} & \thead{\color{ctpBlue}Soal} & \thead{\color{ctpBlue}Poin} & \thead{\color{ctpBlue}Porsi} \\
  \hline
  C2 (memahami) & A1, A4 & 30 & 30\% \\
  \hline
  C3 (menerapkan) & A2 & 15 & 15\% \\
  \hline
  C4 (menganalisis) & A3, B1 & 35 & 35\% \\
  \hline
  C5 (mengevaluasi) & B2 & 20 & 20\% \\
  \hline
  \multicolumn{2}{|r|}{\bfseries Total} & \bfseries 100 & \bfseries 100\% \\
  \hline
\end{tabular}
\end{center}

\vspace{0.2cm}
Estimasi alokasi waktu pengerjaan: A1 \textasciitilde12 menit, A2
\textasciitilde15 menit, A3 \textasciitilde15 menit, A4 \textasciitilde13 menit,
B1 \textasciitilde20 menit, B2 \textasciitilde25 menit (total
\textasciitilde100 menit).

\begin{minipage}{\linewidth}
\begin{bclogo}[couleur=ctpMantle, arrondi=0.1, logo=\bcinfo, ombre=false, couleurBord=ctpSurface, epBord=0.6]{Catatan penskoran}
  Setiap butir dinilai dengan rubrik berjenjang 0--3 pada tiap aspek penilaian.
  Rincian kunci jawaban dan rubrik lengkap terdapat pada dokumen
  \emph{Kunci Jawaban \& Rubrik UTS}, yang bersifat \textbf{internal} dan tidak
  dibagikan kepada mahasiswa.
\end{bclogo}
\end{minipage}

% =============================================================
\section{Panduan Belajar}\label{sec:panduan}

Fokuskan persiapan pada \emph{pemahaman} dan \emph{analisis}, bukan hafalan.
UTS menuntut kemampuan menjelaskan alasan di balik pilihan desain.

\begin{filledstarlist}
  \item Kuasai kerangka \textbf{tiga pilar} dan mampu memetakan setiap tugas
        multimodal ke dalamnya.
  \item Pahami \textbf{intuisi} di balik tiap representasi: mengapa CNN untuk
        gambar, mengapa MFCC/spektrogram untuk audio, mengapa transformer untuk
        konteks teks.
  \item Hafalkan \textbf{trade-off} keempat strategi fusi dan kapan masing-masing
        lebih tepat --- bukan sekadar urutan namanya.
  \item Pahami \textbf{mengapa penyejajaran itu sulit} dan bagaimana attention
        menyelesaikannya secara adaptif.
  \item Latih membaca dan menggambar ulang \textbf{diagram arsitektur} ---
        keterampilan ini langsung diuji pada Butir B1.
  \item Tinjau kembali \textbf{hasil hands-on} Pertemuan 02, 03, 04, dan 06;
        soal dapat menyentuh keputusan implementasi nyata.
\end{filledstarlist}

% =============================================================
\section{Referensi}\label{sec:referensi}

\begin{enumerate}
  \item Baltrušaitis, T., Ahuja, C., \& Morency, L.-P. (2018).
        \emph{Multimodal Machine Learning: A Survey and Taxonomy}.
        IEEE TPAMI, 40(2), 421--443.
  \item Liang, P., Zadeh, A., \& Morency, L.-P. (2023).
        \emph{Tutorial on Multimodal Machine Learning}. ICML.
  \item Goodfellow, I., Bengio, Y., \& Courville, A. (2016).
        \emph{Deep Learning}. MIT Press.
  \item Vaswani, A., et al. (2017). \emph{Attention Is All You Need}. NeurIPS.
  \item Devlin, J., et al. (2019). \emph{BERT: Pre-training of Deep Bidirectional
        Transformers for Language Understanding}. NAACL.
  \item Dosovitskiy, A., et al. (2021). \emph{An Image is Worth 16x16 Words:
        Transformers for Image Recognition at Scale}. ICLR.
  \item Materi kuliah Pertemuan 01--07 beserta lembar kerja (\emph{hands-on})
        terkait.
\end{enumerate}

\end{document}
